\chapter{Implementação do NatureProtector V1}

\section{Visão Geral da Implementação}

Este capítulo descreve a implementação da base técnica de prevenção do NatureProtector V1. A V1 não pretende fornecer um modelo cientificamente calibrado de previsão de incêndios, nem substituir sistemas oficiais de avaliação de perigo ou risco. O objetivo é demonstrar uma pipeline técnica e rastreável, capaz de transformar observações simuladas em avaliações candidatas de risco, projeções operacionais, respostas de API e evidência em tempo de execução. Esta distinção é necessária porque sistemas e índices reconhecidos, como o FWI, produtos do IPMA e plataformas europeias de monitorização, assentam em metodologias e processos de validação próprios~\citep{nrcan_fwi_system,ipma_fwi,ipma_pir_rcm,ec_jrc_effis}.

O fluxo implementado pode ser resumido da seguinte forma:

\begin{center}
\texttt{SensorReadingProduced}
$\rightarrow$
\texttt{RabbitMQ}
$\rightarrow$
\texttt{event\_inbox}
$\rightarrow$
\texttt{NormalizedReading}
$\rightarrow$
\texttt{RiskInput}
$\rightarrow$
\texttt{RiskAssessment}
$\rightarrow$
\texttt{Operational Projections}
$\rightarrow$
\texttt{API / Grafana}
\end{center}

Esta organização separa transporte, validação, avaliação de risco, projeção e observabilidade. Essa separação é essencial para a auditabilidade da V1: um evento recebido pode ser rastreado desde a inbox até às tentativas de processamento, leituras aceites, avaliações de risco, projeções operacionais e, quando aplicável, estado de alerta.

\section{Contratos, Eventos e Pipeline Durável}

A fronteira entre o simulador e o módulo de prevenção é baseada em eventos. O simulador publica leituras de sensores e o \texttt{Prevention.Host} consome essas mensagens a partir do RabbitMQ. Esta separação entre produtor, broker e consumidor é coerente com modelos assíncronos de mensageria, nos quais entrega, consumo e confirmação devem ser tratados como etapas distintas~\citep{rabbitmq_amqp_model,rabbitmq_consumer_acknowledgements}.

O envelope de transporte é representado por \texttt{EventEnvelope<TPayload>}. A sua função é transportar o payload e os metadados necessários à rastreabilidade, como identificador do evento, identificador de correlação, produtor, tipo de evento, versão do esquema, área e instante do evento.

O evento consumido pela V1 é \texttt{SensorReadingProduced}. O respetivo payload contém a leitura simulada e o contexto necessário para a pipeline, incluindo informação sobre sensor, métrica, valor, unidade, localização, instante e execução de simulação.

É importante distinguir contratos de transporte de conceitos internos. \texttt{EventEnvelope<TPayload>} e \texttt{SensorReadingProduced} pertencem à fronteira RabbitMQ. Conceitos como \texttt{OperationalEvent}, \texttt{NormalizedReading}, \texttt{RiskInput} e \texttt{RiskAssessment} pertencem ao processamento interno. Esta separação reduz acoplamento e permite evoluir a pipeline sem alterar necessariamente a fronteira externa de eventos.

Uma decisão central da V1 é não tratar o consumo de mensagens como sinónimo de processamento bem-sucedido. Depois de ser consumido, o evento é persistido em \texttt{pipeline.event\_inbox}. As tentativas são registadas em \texttt{pipeline.processing\_attempts}. Eventos inválidos são classificados em \texttt{pipeline.rejected\_events}. Eventos que esgotam a política de retry ou atingem estado não recuperável são registados em \texttt{pipeline.quarantined\_events}. Esta abordagem aproxima-se de padrões usados em sistemas distribuídos para preservar rastreabilidade e idempotência entre receção da mensagem, tentativa de processamento e efeito persistido~\citep{richardson_idempotent_consumer,richardson_transactional_outbox}.

Esta estrutura permite distinguir quatro situações: evento recebido mas ainda não processado, evento processado com sucesso, evento rejeitado por validação e evento colocado em quarentena por falha ou retry esgotado. Esta distinção é essencial para interpretar corretamente a evidência técnica.

\section{Validação, Elegibilidade e Avaliação de Risco}

Após a ingestão, os eventos são transformados em conceitos internos da pipeline. O objetivo é evitar que a avaliação de risco dependa diretamente de payloads brutos. A pipeline introduz fronteiras de validação, normalização, classificação e elegibilidade.

A normalização produz uma \texttt{NormalizedReading}, isto é, uma leitura numa forma adequada ao processamento posterior. A qualidade e a classificação são representadas através de conceitos como \texttt{QualityFlags} e \texttt{ClassifierResult}. A decisão sobre a utilização da leitura é representada por \texttt{RiskEligibilityResult}.

A V1 distingue três casos principais: leituras completas e elegíveis, leituras parciais mas utilizáveis e leituras bloqueadas. O ponto metodológico mais importante é que uma leitura bloqueada não representa risco baixo. Representa ausência de condições para produzir uma avaliação válida a partir dessa leitura. Esta decisão evita que dados inválidos, incompletos ou inutilizáveis sejam transformados indevidamente em condições seguras.

O \texttt{RiskInput} representa a fronteira entre observações processadas e avaliação de risco. Deve conter a informação necessária para o cálculo, mas não deve conter resultados posteriores, como nível final de risco, estado de alerta ou projeção operacional.

O serviço de cálculo de risco, representado por uma abstração como \texttt{IRiskScoringService}, recebe um \texttt{RiskInput} válido e produz um \texttt{RiskAssessment}. Este resultado representa uma avaliação candidata de risco operacional.

Na V1, esta avaliação deve ser interpretada com prudência. A implementação pode calcular e persistir scores, níveis de risco e resumos, mas estes valores não são medições cientificamente calibradas de perigo de incêndio. São parâmetros e outputs candidatos, usados para exercitar e validar tecnicamente a pipeline. Esta cautela é especialmente importante porque índices e sistemas operacionais de perigo de incêndio, como o FWI ou produtos nacionais/europeus, dependem de entradas, pressupostos e calibração próprios~\citep{nrcan_fwi_system,ipma_fwi,ec_jrc_effis}.

\section{Persistência, Projeções, API e Alertas}

O PostgreSQL é a principal fonte durável da V1. A base de dados organiza-se em três schemas: \texttt{control}, para configuração, áreas, sensores, cenários e execuções de simulação; \texttt{pipeline}, para inbox, tentativas, rejeições e quarentenas; e \texttt{projection}, para leituras aceites, avaliações de risco, snapshots, estado operacional e alertas.

As avaliações de risco não são o output final do sistema. São resultados intermédios que alimentam projeções operacionais. Estas projeções permitem que a API e os painéis de observabilidade exponham estado sem recalcular toda a pipeline.

\begin{table}[h]
\centering
\caption{Principais estruturas persistidas da V1.}
\label{tab:v1-persisted-structures}
\begin{tabular}{p{0.34\textwidth} p{0.52\textwidth}}
\hline
\textbf{Estrutura} & \textbf{Finalidade} \\
\hline
\texttt{control.*} & Configuração da área piloto, sensores, células, cenários e execuções de simulação. \\
\texttt{pipeline.event\_inbox} & Eventos recebidos e respetivo estado de processamento. \\
\texttt{pipeline.processing\_attempts} & Tentativas de processamento, resultados e erros. \\
\texttt{pipeline.rejected\_events} & Eventos rejeitados antes da pipeline de risco aceite. \\
\texttt{pipeline.quarantined\_events} & Eventos com retry esgotado ou estado não recuperável. \\
\texttt{projection.accepted\_reading\_log} & Leituras aceites na pipeline de risco. \\
\texttt{projection.risk\_assessment\_log} & Avaliações candidatas de risco. \\
\texttt{projection.cell\_operational\_state} & Estado operacional por célula. \\
\texttt{projection.area\_operational\_state} & Estado operacional agregado por área. \\
\texttt{projection.alert\_state} & Estado de alerta produzido pela política interna. \\
\hline
\end{tabular}
\end{table}

A Backoffice API expõe o estado operacional produzido pela pipeline. O seu papel é fornecer acesso de leitura a projeções persistidas, como estado operacional da área, score agregado, nível de risco, severidade, resumo, número de avaliações e alertas ativos quando existirem.

A API deve ser interpretada como camada de leitura sobre projeções, não como componente independente de cálculo de risco. Esta decisão evita divergências entre valores persistidos e valores expostos.

A V1 inclui ainda uma política interna de alertas com estados como \texttt{None}, \texttt{Warning} e \texttt{Alarm}. Esta política usa limiares candidatos e histerese para evitar oscilações excessivas quando o score varia perto de um limite. Os limiares usados na V1 não são oficiais nem cientificamente calibrados. Servem para exercitar a pipeline, produzir estado operacional e validar transições através de testes automatizados. Por isso, não devem ser confundidos com classes ou níveis oficiais de perigo de incêndio~\citep{ipma_fwi,ipma_pir_rcm,copernicus_effis}.

A evidência em tempo de execução pode não mostrar alertas ativos se o score agregado estiver abaixo do limiar configurado. Isso não invalida a política. Nesses casos, o comportamento de abertura, descida e resolução deve ser demonstrado por testes.

\section{Observabilidade e Limitações Técnicas}

A observabilidade é apoiada por InfluxDB e Grafana. O InfluxDB suporta dados temporais ou medições de observabilidade, enquanto o Grafana permite construir painéis para inspeção e demonstração. Exemplos de painéis úteis incluem saúde da pipeline, eventos rejeitados e em quarentena, avaliações de risco, estado operacional por área e estado de alertas.

Esta camada não substitui a evidência durável. O PostgreSQL continua a ser a fonte principal para auditoria técnica da pipeline. O Grafana e o InfluxDB apoiam visibilidade, diagnóstico e demonstração, mas não validam cientificamente o modelo.

A implementação V1 fornece uma base funcional e auditável, mas tem limitações relevantes: baseia-se em dados simulados; usa pesos, limiares e classes candidatos; alguns conceitos internos, como \texttt{BaseRisk} ou \texttt{AdjustedScore}, podem estar visíveis na lógica de domínio ou nos testes sem aparecerem como colunas persistidas dedicadas; a política de alertas não equivale a um sistema operacional completo de notificação, escalonamento ou resposta; e a V1 é uma base local, não uma implantação pronta para produção.
Estas limitações não invalidam a V1 enquanto base técnica. Definem os limites da sua interpretação e indicam o trabalho necessário para futuras fases de calibração, validação com dados reais e robustecimento operacional.